% SOURCE_AUTHORITY: sga5_fr_workpass.tex, lines 13660--13693 (LF-normalized unit).
% SOURCE_SHA256: 791F4EFFC5E02832D5D77ED03518C8156D6F07E4C8238B03545DB93D883FBB28
\medskip

\begin{prooftext}[Demostración de la proposición 4.6]
Ya se ha observado que se puede reducir al caso $g = e$, $\varphi = \id_G$. Si se denota por $G_1$ el núcleo del homomorfismo
\[
G \longrightarrow \Aut_k(\fm/\fm^2) \quad ,
\]
entonces, gracias a la teoría de la ramificación, se tiene:
\[
\card(G_1) = p^t \quad .
\]

Por consiguiente, utilizando las notaciones del lema 4.7, se tiene:
\[
\card(G) = p^t c \quad .
\]

Como, por otra parte, $c^{\id_G}(e) = \card(G)$, el enunciado de la proposición 4.6 en el caso $g = e$, $\varphi = \id_G$ y, por consiguiente, en general, resulta del lema 4.7.
\end{prooftext}

\medskip
\noindent
\textbf{Corolario 4.8.} \emph{Para cada $\ell \neq \car(k)$, la imagen del número racional
\[
\frac{(\nu(g_R^{-1}f'_R) - 1)}{c^\varphi(g)}
\]
por la inclusión canónica
\[
\bQ \hookrightarrow \bQ_\ell
\]
es un entero $\ell$-ádico.}

\medskip
\noindent
\textbf{Proposición 4.9.} \emph{La conjetura 4.5 es verdadera si la dimensión del esquema $X$ es igual a 1.}
